\section{Conclusion}
\label{sec:conclusion}

This companion has walked the machinery of Vibe Theory from the bottom up. The substrate is a
growing hyperbolic crystal, a regular Coxeter honeycomb, and we treated it in the detail it
deserves: the reflection groups and Schl\"afli symbols, the regular tilings of the hyperbolic
plane, the heptagrid $\{7,3\}$, the regular honeycombs of hyperbolic space, the dodecagrid
$\{5,3,4\}$ that is the model's real three-dimensional substrate, the classification that pins the
spatial dimension to a window and then to three, the curvature that buys Lorentz invariance, and
the deterministic automaton that grows the crystal one cell at a time forever. On that substrate
the ternary tones and one local rule give an emergent geometry and time, a bounded-below
Hamiltonian, the field content and the gauge and graviton operators from one action, the Newtonian
limit and a conserving Einstein equation, a discrete graviton, black-hole entropy and Hawking
radiation, the quantum pillars with a derived Born rule, a cosmology with an arrow and a dark
sector, sharp predictions that pass the bounds while excluding a lattice, anomaly-forced charge
quantization, baryogenesis, the mass hierarchy, and the measurable structure of the interior.

What the program shows is that a strict monism on an experiential, geometric base is not empty: the
foundation can be made non-arbitrary, the geometry and dimension forced rather than chosen, and a
wide span of physics and a structural account of mind made to come out of one rule, in a finite
simulation, with results that could have failed and sometimes did before being corrected. That is a
reason to take the picture seriously as a hypothesis, not a reason to take it as established.

We have been explicit about the limits. The number of generations and the gauge group are assumed,
not derived. The chiral gauge theory, structure formation, and the dark-sector numbers are first
rungs. And the deepest question, whether the structure of experience the model measures is the same
as the felt fact of experiencing, is a boundary a simulation cannot cross from outside. The honest
bottom line is the one the scoreboard records: a large and growing body of the model is
demonstrated, a real frontier is partial, and a few deep things are open. This is a simulable model
with a concrete build target, not a finished empirical theory, and the geometry at its centre, the
hyperbolic Coxeter crystal grown forever, is the part on which everything else stands.
